
% définir le français comme langue
%\usepackage[french]{babel}\frenchsetup{og=«, fg=»}
\usepackage[autolanguage]{numprint} % for the \nombre command
% pour signaler les erreurs de typo
%\usepackage[All]{lua-typo}
%Hyphenation rules
\usepackage{hyphenat}

\usepackage{tikz,fontspec}
\usepackage{tikz-3dplot}

% fonte pour LuaLaTeX
\usepackage{fontspec}

% ajout des chiffres oldstyle
%\addfontfeature{Numbers={OldStyle,Proportional}}
% pour les formules mathématiques
\usepackage{unicode-math}
\usepackage{amsmath}
%\setmathfont{STIX Two Math}
%\setmainfont{STIX Two Math}
%\usepackage{fancyhdr} 

\usepackage{multicol}

% color
\usepackage{xcolor}
\definecolor{pkcolor}{rgb}{0.0, 0.2, 0.4}  % Blue

\usepackage{titlesec}
\usepackage{titling}
\usepackage{indentfirst}
\usepackage{tabularray}

\usepackage{pgfplots}
\pgfplotsset{compat=1.18}

\usepackage[version=4]{mhchem}
%\usepackage{footnotehyper}
\usepackage{bm}
\usepackage[colorlinks]{hyperref}
\hypersetup{
	colorlinks=true,
	linkcolor=blue,
	filecolor=magenta,      
	urlcolor=cyan,
	citecolor=violet,
	pdftitle={},
	pdfpagemode=FullScreen,
}
\usepackage{afterpage}


% shortcuts
\newcommand{\PQ}{plénodynamique quantique}
\newcommand{\MQ}{mécanique quantique}
\usepackage{xfrac}
%\usepackage{units}
% belles fractions
\newcommand{\R}[2]{\sfrac{#1}{#2}}
\newcommand{\Rv}[4]{\raisebox{0.3ex}{${\scriptscriptstyle #1}$}$\frac{#2}{#3}#4$}
\newcommand{\Ru}[4]{\raisebox{0.7ex}{${\scriptscriptstyle #1}$}$\sfrac{#2}{#3}#4$}


% headers
\usepackage{fancyhdr}

% Redéfinition du style du numéro de page pour les headers
\newcommand{\thepagefhf}{%
	\begingroup
	\setlength{\fboxsep}{3pt} % Espace entre le chiffre et le bord du carré
	\colorbox{pkcolor}{\color{white}\sffamily\small\bfseries\arabic{page}}%
	\endgroup
}

\newcommand{\titlehead}{
	\color{pkcolor}\sffamily\small\bfseries{Plénodynamique quantique}
}

% Redéfinition du style 'plain' (utilisé sur la 1ère page) pour qu'il soit vide
\fancypagestyle{plain}{
	\fancyhf{} % Supprime tout (en-tête et pied de page)
	\renewcommand{\headrulewidth}{0pt} % Enlève la ligne si présente
}

\pagestyle{fancy}
\fancyhf{}
\fancyhead[RO,LE]{\titlehead}
\fancyhead[LO,RE]{\thepagefhf}
% enlève le pied de page
\cfoot{}
% trait horizontale
\renewcommand{\headrulewidth}{1pt}
\renewcommand{\headrule}{\hbox to\headwidth{\color{pkcolor}\leaders\hrule height \headrulewidth\hfill}}



% Title and subtitle
\newcommand{\titlefont}{\bfseries\color{pkcolor}\fontsize{20}{24}\sffamily\selectfont}


% Redefine headers title
\makeatletter

\newcommand{\@subtitle}{}
\newcommand{\subtitle}[1]{\renewcommand{\@subtitle}{#1}}
\newcommand{\subtitlefont}{\bfseries\color{pkcolor}\fontsize{16}{18}\sffamily\selectfont}

\def\@maketitle{
	\begingroup
		{\raggedright\titlefont\@title\par\vskip8pt}
		{\raggedright\subtitlefont\@subtitle\par\vskip8pt}
		{\normalsize\sffamily\raggedright\@author\hfill{Version\,: \Version}\par}
		{\raggedright\sffamily\selectfont\@date\par\vskip18pt}
	\endgroup
}
\makeatother

% redefine section
\setcounter{secnumdepth}{5}

\titleformat{\section}[hang]
	{\color{pkcolor}\sffamily\large\bfseries}
	{\thesection.}
	{0.5em}
	{\uppercase}

\titleformat{name=\section,numberless}[hang]
	{\color{pkcolor}\sffamily\large\bfseries}
	{}
	{0em}
	{\tikz\draw[pkcolor, fill=pkcolor] (0,0) rectangle (1.25ex, 1.25ex); \hspace{2.5pt}\uppercase}

\titleformat{\subsection}[hang]
	{\color{pkcolor}\sffamily\bfseries}
	{\thesubsection.}
	{0.5em}
	{}     

\titleformat{\subsubsection}[block] % You may change it to "runin"
	{\color{pkcolor}\small\sffamily\bfseries\itshape}
	{\thesubsubsection.}
	{0.5em}
	{} % If using runin, change #1 to '#1. ' (space at the end)    

\titleformat{\paragraph}[runin]
	{\small\bfseries}
	{}
	{0em}
	{} 


% citation

% mise en exergue sympa
\usepackage{awesomebox}
\usepackage[most]{tcolorbox}
\usepackage{varwidth}
\tcbuselibrary{most}
\newtcolorbox{mybox}[2][]{enhanced,
	before skip=2mm,after skip=2mm,
	colback=black!5,colframe=black!50,boxrule=0.2mm,
	attach boxed title to top left={xshift=1cm,yshift*=1mm-\tcboxedtitleheight},
	varwidth boxed title*=-3cm,
	boxed title style={frame code={
			\path[fill=tcbcolback!30!black]
			([yshift=-1mm,xshift=-1mm]frame.north west)
			arc[start angle=0,end angle=180,radius=1mm]
			([yshift=-1mm,xshift=1mm]frame.north east)
			arc[start angle=180,end angle=0,radius=1mm];
			\path[left color=tcbcolback!60!black,right color=tcbcolback!60!black,
			middle color=tcbcolback!80!black]
			([xshift=-2mm]frame.north west) -- ([xshift=2mm]frame.north east)
			[rounded corners=1mm]-- ([xshift=1mm,yshift=-1mm]frame.north east)
			-- (frame.south east) -- (frame.south west)
			-- ([xshift=-1mm,yshift=-1mm]frame.north west)
			[sharp corners]-- cycle;
		},interior engine=empty,
	},
	fonttitle=\bfseries,
	title={#2},#1}

% symboles
\usepackage{pifont}% http://ctan.org/pkg/pifont
\usepackage[dvipsnames]{xcolor}
\newcommand{\vrai}[1][black]{\textcolor{#1}{\ding{51}}}
\newcommand{\faux}[1][black]{\textcolor{#1}{\ding{55}}}
\newcommand{\vraiV}{\vrai[OliveGreen]}%
\newcommand{\fauxR}{\faux[red]}%
% figure

\usepackage[section]{placeins}
\usepackage{caption, subcaption}
% same as tabular environment
\captionsetup[table]{
	position=top,           % ← c'est la ligne importante
%	skip=10pt,              % espace entre le caption et le tableau
%	labelfont=bf,           % label en gras (Table 1)
	textfont=normal         % texte du caption en normal
}

\usepackage{floatrow}
\usepackage{float}
\usepackage{wrapfig}
\usepackage{graphicx}
\graphicspath{ {./figures/} }


% biblio
\usepackage{csquotes}
\usepackage[
backend=biber,
%style=phys,
sorting=ynt,
citestyle=numeric-comp,
]{biblatex}

\addbibresource{plenodynamique-quantique.bib}

% Nœud
\newtcbox{\vnode}{%
	boxrule=0.4pt,          
	arc=1pt,                
	colback=gray!15,        
	colframe=gray!60,        
	coltext=pkcolor,         
	left=1pt, right=1pt,    
	top=1pt, bottom=1pt,      
	boxsep=0pt,             
	fontupper=\ttfamily\small,
	valign=center,
	nobeforeafter,
	box align=base
}
\newtcbox{\node}{%
	boxrule=0.4pt,          
	arc=1pt,                
	colback=blue!15,        
	colframe=blue!60,        
	coltext=black,         
	left=1pt, right=1pt,    
	top=1pt, bottom=1pt,      
	boxsep=0pt,             
	fontupper=\ttfamily\small,
	valign=center,
	nobeforeafter,
	box align=base
}

%\newcommand{\vvnode}[2]{\vnode{$\rotatebox{180}{S}_#1}{#2}}$}

% Lists

\usepackage{enumitem}

\newcommand{\squarebullet}[1][pkcolor]{%
	\tikz[baseline=0.0ex]{
		\fill[#1] (0,0) rectangle (1.1ex,1.1ex);
	}%
}

\setlist[itemize]{
	%label=\textbullet,        % puce classique
	 label=\squarebullet[pkcolor],   % exemple de puce colorée
}